% chapters/16/section-02.tex
\section{Checklist Fungsional}
\label{sec:checklist-fungsional}

\subsection{Menurunkannya dari GDD}

\istilah{Checklist} fungsional adalah daftar pernyataan yang harus bernilai
benar sebelum \istilah{build} dianggap selesai. Ia bukan daftar keinginan,
melainkan daftar janji: setiap butirnya berasal dari sesuatu yang sudah Anda
tuliskan pada \istilah{Game Design Document} (GDD) Anda.

Cara menyusunnya tidak menuntut pekerjaan baru, hanya ketelitian membaca
kembali. Buka GDD Anda, lalu ubah setiap pernyataan menjadi satu butir yang
dapat diperiksa benar atau salah. Tiga contoh penurunannya diperlihatkan pada
Tabel~\ref{tab:turun-checklist}.

\begin{table}[htbp]
\centering
\caption{Menurunkan butir \istilah{checklist} dari pernyataan GDD}
\label{tab:turun-checklist}
{\setlength{\tabcolsep}{4pt}\footnotesize%
\begin{tabularx}{\textwidth}{@{}L{5.2cm} Y@{}}
\toprule
\headrow \textbf{Pernyataan pada GDD} & \textbf{Butir \istilah{checklist} yang dapat diperiksa} \\
\midrule
``Pemain mengumpulkan kunci untuk membuka pintu akhir.'' &
Benda bertanda \code{Pickup} yang tersedia berjumlah sekurang-kurangnya tiga,
dan pintu akhir benar-benar terbuka setelah ketiganya diambil. \\
\addlinespace
``Permainan dapat diselesaikan dalam lima sampai sepuluh menit.'' &
Satu kali penyelesaian penuh berlangsung antara lima dan sepuluh menit pada
permainan yang dijalankan pertama kali. \\
\addlinespace
``Pemain yang jatuh tidak langsung kehilangan seluruh kemajuannya.'' &
Tidak ada tempat di level yang membuat pemain tidak dapat melanjutkan
permainan, dan tidak ada kemajuan yang hilang lebih dari satu bagian. \\
\bottomrule
\end{tabularx}}
\end{table}

Perhatikan butir ketiga. Pernyataan GDD-nya bernilai rancangan, sedangkan
butir \istilah{checklist}-nya bernilai pemeriksaan, dan di situlah letak
gunanya: ia mengubah niat menjadi syarat yang dapat diuji, termasuk oleh
pemeriksaan otomatis pada Bagian~\ref{sec:uji-otomatis}.

\subsection{Tiga Kelompok Butir}

\istilah{Checklist} yang hanya memuat satu jenis butir akan meninggalkan
bagian yang penting. Susunlah butirnya dalam tiga kelompok berikut.

\begin{enumerate}
  \item \textbf{Mekanik.} Apakah setiap mekanik yang dijanjikan GDD benar-benar
        ada dan berjalan? Lompatan, pengambilan benda, kondisi menang,
        kondisi kalah.
  \item \textbf{Integrasi.} Apakah bagian-bagian yang dibuat pada bab yang
        berbeda benar-benar tersambung? Bunyi terdengar pada peristiwanya,
        animasi berganti menurut keadaannya, antarmuka menampilkan angka yang
        benar.
  \item \textbf{Keseluruhan.} Apakah permainannya dapat diselesaikan dari awal
        sampai akhir, tanpa berhenti, tanpa jalan buntu, dan dalam waktu yang
        Anda janjikan?
\end{enumerate}

Kelompok ketiga adalah yang paling sering terlewat, sebab ia tidak dapat
diperiksa sambil membuat. Ia hanya dapat diperiksa dengan menjalankan
permainannya dari awal sampai akhir, dan itulah sebabnya ia memerlukan
\istilah{build}, bukan \istilah{scene} yang sedang Anda sunting. Kebiasaan
menyusun daftar periksa seperti ini juga dipakai pada penyusunan kurikulum dan
rencana pembelajaran bidang permainan \cite{polibatam-kurikulum}.

\subsection{Checklist Fungsional sebagai Berkas}

\istilah{Checklist} Anda sebaiknya berupa berkas, bukan hafalan atau ingatan.
Berkas itu memudahkan Anda menjalankannya berulang kali dan mencatat hasilnya
setiap kali, sehingga kemajuan antara satu putaran dan putaran berikutnya
terlihat. Bentuknya diperlihatkan pada Tabel~\ref{tab:checklist-kosong}.

\begin{table}[htbp]
\centering
\caption{Bentuk \istilah{checklist} fungsional yang diisi setiap putaran pengujian}
\label{tab:checklist-kosong}
{\setlength{\tabcolsep}{4pt}\footnotesize%
\begin{tabularx}{\textwidth}{@{}C{0.8cm} L{3.0cm} Y C{1.6cm}@{}}
\toprule
\headrow \textbf{No} & \textbf{Kelompok} & \textbf{Butir pemeriksaan} & \textbf{Putaran 1 / 2 / 3} \\
\midrule
1 & Mekanik & & \\[1.3em]
\midrule
2 & Mekanik & & \\[1.3em]
\midrule
3 & Integrasi & & \\[1.3em]
\midrule
4 & Integrasi & & \\[1.3em]
\midrule
5 & Keseluruhan & & \\[1.3em]
\midrule
\addlinespace
\multicolumn{4}{@{}p{\textwidth}@{}}{\footnotesize\itshape
Jumlah butir yang lulus: \dotfill} \\[1.4em]
\addlinespace
\multicolumn{4}{@{}p{\textwidth}@{}}{\footnotesize\itshape
Butir yang belum lulus dan alasan belum diperbaikinya: \dotfill} \\[1.8em]
\bottomrule
\end{tabularx}}
\end{table}

\begin{tipsbox}
Batasi \istilah{checklist} Anda pada sekitar dua puluh sampai tiga puluh butir.
Daftar yang terlalu panjang tidak akan pernah dijalankan sampai selesai,
sedangkan daftar yang terlalu pendek akan melewatkan hal yang penting. Butir
yang paling berguna adalah butir yang pernah gagal pada putaran sebelumnya.
\end{tipsbox}
